\documentclass[12pt,letterpaper]{article}
\usepackage{graphicx,textcomp}
\usepackage{natbib}
\usepackage{setspace}
\usepackage{fullpage}
\usepackage{color}
\usepackage[reqno]{amsmath}
\usepackage{amsthm}
\usepackage{fancyvrb}
\usepackage{amssymb,enumerate}
\usepackage[all]{xy}
\usepackage{endnotes}
\usepackage{lscape}
\newtheorem{com}{Comment}
\usepackage{float}
\usepackage{hyperref}
\newtheorem{lem} {Lemma}
\newtheorem{prop}{Proposition}
\newtheorem{thm}{Theorem}
\newtheorem{defn}{Definition}
\newtheorem{cor}{Corollary}
\newtheorem{obs}{Observation}
\usepackage[compact]{titlesec}
\usepackage{dcolumn}
\usepackage{tikz}
\usetikzlibrary{arrows}
\usepackage{multirow}
\usepackage{xcolor}
\newcolumntype{.}{D{.}{.}{-1}}
\newcolumntype{d}[1]{D{.}{.}{#1}}
\definecolor{light-gray}{gray}{0.65}
\usepackage{url}
\usepackage{listings}

\definecolor{codegreen}{rgb}{0,0.6,0}
\definecolor{codegray}{rgb}{0.5,0.5,0.5}
\definecolor{codepurple}{rgb}{0.58,0,0.82}
\definecolor{backcolour}{rgb}{0.95,0.95,0.92}

\lstdefinestyle{mystyle}{
	backgroundcolor=\color{backcolour},   
	commentstyle=\color{codegreen},
	keywordstyle=\color{magenta},
	numberstyle=\tiny\color{codegray},
	stringstyle=\color{codepurple},
	basicstyle=\footnotesize,
	breakatwhitespace=false,
	breaklines=true,
	captionpos=b,
	keepspaces=true,
	numbers=left,
	numbersep=5pt,
	showspaces=false,
	showstringspaces=false,
	showtabs=false,
	tabsize=2
}
\lstset{style=mystyle}

\newcommand{\Sref}[1]{Section~\ref{#1}}
\newtheorem{hyp}{Hypothesis}

\title{Problem Set 3}
\date{Due: March 25, 2026}
\author{Applied Stats II}

\begin{document}
	\maketitle
	
	\section*{Instructions}
	\begin{itemize}
		\item Please show your work! You may lose points by simply writing in the answer. If the problem requires you to execute commands in \texttt{R}, please include the code you used to get your answers. Please also include the \texttt{.R} file that contains your code. If you are not sure if work needs to be shown for a particular problem, please ask.
		\item Your homework should be submitted electronically on GitHub in \texttt{.pdf} form.
		\item This problem set is due before 23:59 on Wednesday March 25, 2026. No late assignments will be accepted.
	\end{itemize}
	
	\vspace{.25cm}
	\section*{Question 1}
	\vspace{.25cm}
	\noindent We are interested in how governments' management of public resources impacts economic prosperity. Our data come from \href{https://www.researchgate.net/profile/Adam_Przeworski/publication/240357392_Classifying_Political_Regimes/links/0deec532194849aefa000000/Classifying-Political-Regimes.pdf}{Alvarez, Cheibub, Limongi, and Przeworski (1996)} and is labelled \texttt{gdpChange.csv} on GitHub. The dataset covers 135 countries observed between 1950 or the year of independence or the first year for which data on economic growth are available ("entry year"), and 1990 or the last year for which data on economic growth are available ("exit year"). The unit of analysis is a particular country during a particular year, for a total $>$ 3,500 observations. 
	
	\begin{itemize}
		\item Response variable: 
		\begin{itemize}
			\item \texttt{GDPWdiff}: Difference in GDP between year $t$ and $t-1$. Possible categories include: "positive", "negative", or "no change"
		\end{itemize}
		\item Explanatory variables: 
		\begin{itemize}
			\item \texttt{REG}: 1=Democracy; 0=Non-Democracy
			\item \texttt{OIL}: 1=if the average ratio of fuel exports to total exports in 1984-86 exceeded 50\%; 0= otherwise
		\end{itemize}
	\end{itemize}
	
	\newpage
	\noindent Please answer the following questions:
	
	\begin{enumerate}
		\item Construct and interpret an unordered multinomial logit with \texttt{GDPWdiff} as the output and "no change" as the reference category, including the estimated cutoff points and coefficients.
		
		\lstinputlisting[language=R, firstline=48, lastline=56]{PS3.R}
		
		\textit{The unordered multinomial logit model estimates the effect of regime type (REG) and oil dependence (OIL) on the likelihood of experiencing negative or positive GDP change relative to the reference category of no change. 
			The coefficient for REG is positive for both negative (1.38) and positive (1.77) GDP changes, suggesting that democracies have higher log-odds of experiencing both negative and positive GDP changes relative to no change compared to non-democracies. 
			The estimated standard errors for REG (0.77 for both outcomes) indicate moderate uncertainty around these estimates. The coefficient for OIL is also positive for both negative (4.78) and positive (4.58) GDP changes,
			implying that oil-exporting countries are more likely to experience changes in GDP. However, the standard errors for OIL are very large, suggesting that these estimates are imprecise. 
			Question 1b ordered logit }
			
				\begin{verbatim}
			Coefficients:
			(Intercept)      REG      OIL
			negative    3.805370 1.379282 4.783968
			positive    4.533759 1.769007 4.576321
			
			Std. Errors:
			(Intercept)       REG      OIL
			negative   0.2706832 0.7686958 6.885366
			positive   0.2692006 0.7670366 6.885097
			
			Residual Deviance: 4678.77 
			AIC: 4690.77 
			\end{verbatim}
		
		\item Construct and interpret an ordered multinomial logit with \texttt{GDPWdiff} as the outcome variable, including the estimated cutoff points and coefficients.
		
		\lstinputlisting[language=R, firstline=62, lastline=66]{PS3.R}
		
		\textit{The ordered logit model estimates the effect of regime type (REG) and oil dependence (OIL) 
			on the likelihood of experiencing higher GDP changes, with the outcome ordered from 
			negative to no change to positive. 
			
			The coefficient for REG is 0.3985 (standard error = 0.075), indicating that democracies 
			are more likely than non-democracies to experience higher GDP changes, moving from 
			negative to no change or from no change to positive. 
			
			The coefficient for OIL is -0.1987 (standard error = 0.116), suggesting that 
			oil-exporting countries may be slightly less likely to experience higher GDP changes.} 
			
					\begin{verbatim}
			Coefficients:
			Value Std. Error t value
			REG  0.3985    0.07518   5.300
			OIL -0.1987    0.11572  -1.717
			
			Intercepts:
			Value    Std. Error t value 
			negative|no_change  -0.7312   0.0476   -15.3597
			no_change|positive  -0.7105   0.0475   -14.9554
			
			Residual Deviance: 4687.689 
			AIC: 4695.689  
			\end{verbatim}
	\end{enumerate}
	
	\section*{Question 2} 
	\vspace{.25cm}
	
	\noindent Consider the data set \texttt{MexicoMuniData.csv}, which includes municipal-level information from Mexico. The outcome of interest is the number of times the winning PAN presidential candidate in 2006 (\texttt{PAN.visits.06}) visited a district leading up to the 2009 federal elections, which is a count. Our main predictor of interest is whether the district was highly contested, or whether it was not (the PAN or their opponents have electoral security) in the previous federal elections during 2000 (\texttt{competitive.district}), which is binary (1=close/swing district, 0="safe seat"). We also include \texttt{marginality.06} (a measure of poverty) and \texttt{PAN.governor.06} (a dummy for whether the state has a PAN-affiliated governor) as additional control variables. 
	
	\begin{enumerate}
		\item [(a)] Run a Poisson regression because the outcome is a count variable. Is there evidence that PAN presidential candidates visit swing districts more? Provide a test statistic and p-value.
		
		\textit{ The Poisson regression examines the number of visits by the winning PAN presidential candidate in 2006. 
			There is no evidence that swing districts received more visits than safe districts (competitive.district coefficient = -0.081, z = -0.477, p = 0.634). }
		
		\lstinputlisting[language=R, firstline=82, lastline=85]{PS3.R}
		
					\begin{verbatim}
		Coefficients:
		Estimate Std. Error z value Pr(>|z|)    
		(Intercept)          -3.81023    0.22209 -17.156   <2e-16 ***
		competitive.district -0.08135    0.17069  -0.477   0.6336    
		marginality.06       -2.08014    0.11734 -17.728   <2e-16 ***
		PAN.governor.06      -0.31158    0.16673  -1.869   0.0617 .  
		---
		Signif. codes:  0 ‘***’ 0.001 ‘**’ 0.01 ‘*’ 0.05 ‘.’ 0.1 ‘ ’ 1
		
		(Dispersion parameter for poisson family taken to be 1)
		
		Null deviance: 1473.87  on 2406  degrees of freedom
		Residual deviance:  991.25  on 2403  degrees of freedom
		AIC: 1299.2
		
		Number of Fisher Scoring iterations: 7 
		\end{verbatim}
		
		\item [(b)] Interpret the \texttt{marginality.06} and \texttt{PAN.governor.06} coefficients.
		
		\textit{ Higher poverty is associated with fewer visits by the PAN candidate. In other words, poorer districts tend to receive less attention. PAN.governor.06: Districts in states with a PAN governor tend to get slightly fewer visits (coefficient = -0.312), though the effect is only marginally significant (p ≈ 0.062)}
		
		\item [(c)] Provide the estimated mean number of visits from the winning PAN presidential candidate for a hypothetical district that was competitive (\texttt{competitive.district}=1), had an average poverty level (\texttt{marginality.06} = 0), and a PAN governor (\texttt{PAN.governor.06}=1).
		
		\lstinputlisting[language=R, firstline=94, lastline=100]{PS3.R}
		
					\begin{verbatim}
			> predicted_visits
			1 
			0.01494818 
		\end{verbatim}
		
		\textit{For a competitive district with average poverty and a PAN governor, the predicted mean number of PAN visits is essentially zero (≈0.015). This reflects that most districts received very few visits}
	\end{enumerate}
	
\end{document}